\documentclass[11pt,a4paper]{article}
\usepackage[utf8]{inputenc}
\usepackage[margin=2.5cm]{geometry}
\usepackage{amsmath,amssymb,amsfonts}
\usepackage{array}
\usepackage{hyperref}

\hypersetup{
    colorlinks=true,
    linkcolor=blue,
    urlcolor=blue,
    citecolor=blue
}

\setlength{\parindent}{0pt}
\setlength{\parskip}{6pt}

\newcommand{\kotakdefinisi}[2]{%
\noindent\fbox{%
\parbox{\dimexpr\textwidth-2\fboxsep-2\fboxrule\relax}{%
\textbf{#1}\\[4pt]
#2
}%
}\par\vspace{6pt}
}

\begin{document}

% --- Header Catatan Kuliah ---
\begin{center}
    {\LARGE \textbf{CATATAN KULIAH MA4171 TEORI KONTROL LINEAR}}\\[6pt]
    {\Large \textbf{Topik 03: Transformasi Z dan Sistem Diskrit}}\\[4pt]
    \textit{Cakupan: Minggu ke-5}\\[2pt]
    \textbf{Referensi Utama:} Katsuhiko Ogata, \textit{Modern Control Engineering} \& \textit{Discrete-Time Control Systems}
\end{center}
\vspace{-4pt}
\hrule
\vspace{10pt}

\tableofcontents
\vspace{10pt}
\hrule
\vspace{12pt}

\section{Pendahuluan Sistem Waktu Diskrit \& Pencuplikan}

Dalam implementasi sistem kendali modern, pengontrol umumnya direalisasikan menggunakan komputer digital, mikrokontroler, atau DSP (\textit{Digital Signal Processor}). Sinyal kontinu dari plant fisik harus diubah menjadi deret diskrit melalui proses pencuplikan (\textit{sampling}), dan sebaliknya keluaran komputer diubah kembali menjadi sinyal kontinu menggunakan rangkaian penahan (\textit{hold circuit}).

\subsection{Klasifikasi Sinyal}
\begin{itemize}
    \item \textbf{Sinyal Kontinu (Analog)}: Didefinisikan untuk setiap nilai waktu $t \in \mathbb{R}$ kontinu, dinotasikan $x(t)$.
    \item \textbf{Sinyal Diskrit (Discrete-Time)}: Didefinisikan hanya pada titik-titik waktu diskrit $t = kT$ dengan $k \in \mathbb{Z}$, dinotasikan $x(k)$ atau $x(kT)$, di mana $T$ adalah periode sampling.
    \item \textbf{Sinyal Digital}: Sinyal waktu diskrit yang amplitudonya telah dikuantisasi ke dalam level-level diskrit tertentu (representasi biner terbatas).
\end{itemize}

\subsection{Proses Pencuplikan Ideal \& Teorema Sampling Shannon}
Proses pencuplikan ideal memodelkan sinyal tercuplik $x^*(t)$ sebagai modulasi kereta impuls Dirac:
\begin{equation*}
    x^*(t) = \sum_{k=0}^{\infty} x(kT) \delta(t - kT)
\end{equation*}
Di mana $\delta(t)$ adalah impuls Dirac dan $T$ adalah periode sampling ($f_s = \frac{1}{T}$, $\omega_s = \frac{2\pi}{T}$).

\kotakdefinisi{Teorema Sampling Nyquist-Shannon:}{
Agar sinyal kontinu $x(t)$ dengan batas frekuensi tertinggi $\omega_{max}$ dapat direkonstruksi secara sempurna dari sinyal sampel $x(kT)$, frekuensi pencuplikan $\omega_s$ harus memenuhi:
\begin{equation*}
    \omega_s \ge 2 \omega_{max} \quad \left( \text{atau } f_s \ge 2 f_{max} \right)
\end{equation*}
Jika syarat ini dilanggar, akan timbul fenomena tumpang tindih spektrum frekuensi yang disebut \textbf{aliasing}.
}

\subsection{Penahan Orde Nol (\textit{Zero-Order Hold} / ZOH)}
Untuk menghubungkan sinyal kontrol diskrit $u(kT)$ ke plant fisik kontinu, digunakan penahan orde nol (ZOH) yang mempertahankan nilai sinyal konstan selama interval satu periode sampling $T$:
\begin{equation*}
    u(t) = u(kT) \quad \text{untuk } kT \le t < (k+1)T
\end{equation*}
Fungsi transfer dari ZOH dalam domain Laplace adalah:
\begin{equation*}
    G_h(s) = \frac{1 - e^{-Ts}}{s}
\end{equation*}

\vspace{6pt}

\section{Definisi Transformasi Z}

Transformasi Z merupakan ekuivalen domain diskrit dari Transformasi Laplace pada sistem kontinu.

\kotakdefinisi{Definisi Transformasi Z Satu-Sisi (Unilateral):}{
Diberikan deret waktu diskrit $x(k) = x(kT)$ untuk $k \ge 0$, Transformasi Z didefinisikan sebagai:
\begin{equation*}
    \mathcal{Z}\{x(k)\} = X(z) = \sum_{k=0}^{\infty} x(k) z^{-k}
\end{equation*}
Di mana $z$ adalah variabel kompleks yang berelasi dengan variabel Laplace $s$ melalui:
\begin{equation*}
    z = e^{sT}
\end{equation*}
}

\subsection{Daerah Konvergensi (\textit{Region of Convergence} / ROC)}
Transformasi Z hanya berlaku pada daerah di bidang kompleks $z$ di mana deret tak hingga $\sum_{k=0}^{\infty} x(k) z^{-k}$ konvergen mutlak:
\begin{equation*}
    \sum_{k=0}^{\infty} |x(k) z^{-k}| < \infty
\end{equation*}
Untuk sinyal kausal ($x(k) = 0$ untuk $k < 0$), ROC umumnya berbentuk bagian luar suatu lingkaran pada bidang kompleks: $|z| > R$.

\subsection{Tabel Pasangan Transformasi Z Sinyal Dasar}
Berikut adalah pasangan Transformasi Z untuk sinyal-sinyal dasar yang sering digunakan:
\begin{center}
\small
\begin{tabular}{|l|c|c|c|}
\hline
\textbf{Nama Sinyal} & $x(t), t \ge 0$ & $x(k), k \ge 0$ & $X(z) = \mathcal{Z}\{x(k)\}$ \\ \hline
\textbf{Impuls Satuan} & $\delta(t)$ & $\delta(k) = \begin{cases} 1, & k=0 \\ 0, & k \ne 0 \end{cases}$ & $1$ \\ \hline
\textbf{Tangga Satuan (Step)} & $1(t)$ & $1(k) = 1$ & $\dfrac{z}{z-1} = \dfrac{1}{1-z^{-1}}$ \\ \hline
\textbf{Ramp Satuan} & $t$ & $kT$ & $\dfrac{T z}{(z-1)^2}$ \\ \hline
\textbf{Polinomial $k$} & - & $k$ & $\dfrac{z}{(z-1)^2}$ \\ \hline
\textbf{Eksponensial} & $e^{-at}$ & $e^{-akT} = (e^{-aT})^k$ & $\dfrac{z}{z - e^{-aT}}$ \\ \hline
\textbf{Pangkat Basis $a$} & - & $a^k$ & $\dfrac{z}{z-a}$ \\ \hline
\textbf{Sinusoidal} & $\sin(\omega t)$ & $\sin(\omega k T)$ & $\dfrac{z \sin(\omega T)}{z^2 - 2z \cos(\omega T) + 1}$ \\ \hline
\textbf{Cosinusoidal} & $\cos(\omega t)$ & $\cos(\omega k T)$ & $\dfrac{z (z - \cos(\omega T))}{z^2 - 2z \cos(\omega T) + 1}$ \\ \hline
\end{tabular}
\end{center}

\vspace{6pt}

\section{Sifat-Sifat dan Teorema Transformasi Z}

\begin{enumerate}
    \item \textbf{Linearitas}:\\
    Untuk setiap skalar $\alpha, \beta \in \mathbb{C}$:
    \begin{equation*}
        \mathcal{Z}\{\alpha x_1(k) + \beta x_2(k)\} = \alpha X_1(z) + \beta X_2(z)
    \end{equation*}

    \item \textbf{Teorema Pergeseran Waktu (Time Shifting)}:
    \begin{itemize}
        \item \textbf{Pergeseran Mundur (Delay / Backward Shift)}:\\
        Jika $x(k) = 0$ untuk $k < 0$, maka untuk bilangan bulat $m \ge 0$:
        \begin{equation*}
            \mathcal{Z}\{x(k-m)\} = z^{-m} X(z)
        \end{equation*}
        \item \textbf{Pergeseran Maju (Advance / Forward Shift)}:
        \begin{align*}
            \mathcal{Z}\{x(k+1)\} &= z X(z) - z x(0) \\
            \mathcal{Z}\{x(k+2)\} &= z^2 X(z) - z^2 x(0) - z x(1) \\
            \mathcal{Z}\{x(k+m)\} &= z^m X(z) - \sum_{j=0}^{m-1} z^{m-j} x(j)
        \end{align*}
    \end{itemize}

    \item \textbf{Penskalaan Kompleks}:
    \begin{equation*}
        \mathcal{Z}\{a^{-k} x(k)\} = X(az), \quad \mathcal{Z}\{e^{-akT} x(k)\} = X(z e^{aT})
    \end{equation*}

    \item \textbf{Diferensiasi Kompleks}:
    \begin{equation*}
        \mathcal{Z}\{k \, x(k)\} = -z \frac{d}{dz} X(z)
    \end{equation*}

    \item \textbf{Teorema Nilai Awal (\textit{Initial Value Theorem})}:
    \begin{equation*}
        x(0) = \lim_{z \to \infty} X(z)
    \end{equation*}

    \item \textbf{Teorema Nilai Akhir (\textit{Final Value Theorem})}:\\
    Jika semua kutub dari $(1 - z^{-1})X(z)$ atau $(z-1)X(z)$ berada di dalam lingkaran satuan $|z| < 1$, maka:
    \begin{equation*}
        \lim_{k \to \infty} x(k) = \lim_{z \to 1} (1 - z^{-1}) X(z) = \lim_{z \to 1} (z-1) X(z)
    \end{equation*}

    \item \textbf{Konvolusi Kompleks Diskrit}:\\
    Jika $y(k) = x_1(k) * x_2(k) = \sum_{j=0}^{k} x_1(j) x_2(k-j)$, maka:
    \begin{equation*}
        Y(z) = X_1(z) \cdot X_2(z)
    \end{equation*}
\end{enumerate}

\vspace{6pt}

\section{Invers Transformasi Z}

Untuk mengembalikan fungsi domain $z$, $X(z)$, ke domain waktu diskrit $x(k) = \mathcal{Z}^{-1}\{X(z)\}$, terdapat 3 metode utama:

\subsection{1. Metode Ekspansi Pecahan Parsial (\textit{Partial Fraction Expansion})}
Bentuk yang paling efisien adalah melakukan ekspansi pada fungsi $\dfrac{X(z)}{z}$:
\begin{equation*}
    \frac{X(z)}{z} = \frac{c_1}{z - p_1} + \frac{c_2}{z - p_2} + \dots + \frac{c_n}{z - p_n}
\end{equation*}
Sehingga:
\begin{equation*}
    X(z) = \frac{c_1 z}{z - p_1} + \frac{c_2 z}{z - p_2} + \dots + \frac{c_n z}{z - p_n}
\end{equation*}
Dengan menerapkan inversi langsung dari tabel dasar:
\begin{equation*}
    x(k) = c_1 (p_1)^k + c_2 (p_2)^k + \dots + c_n (p_n)^k, \quad k \ge 0
\end{equation*}

\subsection{2. Metode Pembagian Langsung (\textit{Power Series / Direct Division})}
Fungsi $X(z) = \frac{N(z)}{D(z)}$ dibagi secara langsung dalam pangkat turun $z^{-1}$:
\begin{equation*}
    X(z) = x(0) + x(1) z^{-1} + x(2) z^{-2} + x(3) z^{-3} + \dots
\end{equation*}
Koefisien dari suku $z^{-k}$ secara langsung merupakan nilai sampel $x(k)$.

\subsection{3. Metode Integral Inversi (Teorema Residu)}
\begin{equation*}
    x(k) = \frac{1}{2\pi j} \oint_{C} X(z) z^{k-1} \, dz = \sum_{i} \operatorname{Res}\left[ X(z) z^{k-1}, z = p_i \right]
\end{equation*}
Di mana untuk kutub sederhana $p_i$:
\begin{equation*}
    \operatorname{Res}\left[ X(z) z^{k-1}, z = p_i \right] = \lim_{z \to p_i} (z - p_i) X(z) z^{k-1}
\end{equation*}

\vspace{6pt}

\section{Fungsi Pulsa Transfer (\textit{Pulse Transfer Function})}

\subsection{Persamaan Beda Linear Time-Invariant}
Dinamika sistem diskrit LTI orde-$n$ dimodelkan oleh persamaan beda linear:
\begin{equation*}
    y(k) + a_1 y(k-1) + \dots + a_n y(k-n) = b_0 u(k) + b_1 u(k-1) + \dots + b_m u(k-m)
\end{equation*}
Dengan menerapkan Transformasi Z pada kedua sisi dan mengasumsikan kondisi awal nol:
\begin{equation*}
    (1 + a_1 z^{-1} + \dots + a_n z^{-n}) Y(z) = (b_0 + b_1 z^{-1} + \dots + b_m z^{-m}) U(z)
\end{equation*}

\kotakdefinisi{Fungsi Pulsa Transfer $G(z)$:}{
\begin{equation*}
    G(z) = \frac{Y(z)}{U(z)} = \frac{b_0 + b_1 z^{-1} + \dots + b_m z^{-m}}{1 + a_1 z^{-1} + \dots + a_n z^{-n}} = \frac{b_0 z^n + b_1 z^{n-1} + \dots + b_m z^{n-m}}{z^n + a_1 z^{n-1} + \dots + a_n}
\end{equation*}
}

\subsection{Diskritisasi Plant Kontinu dengan ZOH}
Jika suatu plant kontinu $G_p(s)$ didahului oleh penahan orde nol ZOH $G_h(s) = \frac{1 - e^{-Ts}}{s}$, fungsi transfer diskrit ekuivalennya adalah:
\begin{align*}
    G(z) &= \mathcal{Z}\{ G_h(s) G_p(s) \} = \mathcal{Z}\left\{ \frac{1 - e^{-Ts}}{s} G_p(s) \right\} \\
    &= (1 - z^{-1}) \mathcal{Z}\left\{ \frac{G_p(s)}{s} \right\}
\end{align*}

\kotakdefinisi{Rumus Diskritisasi Plant Kontinu dengan ZOH:}{
\begin{equation*}
    G(z) = (1 - z^{-1}) \mathcal{Z}\left\{ \mathcal{L}^{-1}\left[ \frac{G_p(s)}{s} \right] \right\}
\end{equation*}
}

\vspace{6pt}

\section{Representasi Ruang Keadaan Sistem Diskrit}

\subsection{Bentuk Standar Sistem Ruang Keadaan Diskrit}

\kotakdefinisi{Persamaan Ruang Keadaan LTI Diskrit:}{
\begin{align*}
    x(k+1) &= G x(k) + H u(k) \quad &\text{(\textbf{Persamaan Keadaan Diskrit})} \\
    y(k) &= C x(k) + D u(k) \quad &\text{(\textbf{Persamaan Keluaran Diskrit})}
\end{align*}
}
Di mana:
\begin{itemize}
    \item $x(k) \in \mathbb{R}^n$ : Vektor keadaan pada langkah ke-$k$.
    \item $u(k) \in \mathbb{R}^m$ : Vektor masukan kontrol pada langkah ke-$k$.
    \item $y(k) \in \mathbb{R}^p$ : Vektor keluaran pada langkah ke-$k$.
    \item $G \in \mathbb{R}^{n \times n}$ : Matriks sistem diskrit (beberapa literatur menuliskan $A_d$ atau $\Phi$).
    \item $H \in \mathbb{R}^{n \times m}$ : Matriks masukan diskrit (beberapa literatur menuliskan $B_d$ atau $\Gamma$).
    \item $C \in \mathbb{R}^{p \times n}$ : Matriks keluaran.
    \item $D \in \mathbb{R}^{p \times m}$ : Matriks transmisi langsung.
\end{itemize}

\subsection{Diskritisasi dari Ruang Keadaan Kontinu}
Diberikan sistem kontinu:
\begin{equation*}
    \dot{x}(t) = A x(t) + B u(t), \quad y(t) = C x(t) + D u(t)
\end{equation*}
Dengan masukan dikendalikan oleh ZOH ($u(t) = u(kT)$ untuk $kT \le t < (k+1)T$), solusi kontinu pada interval sampling adalah:
\begin{equation*}
    x((k+1)T) = e^{AT} x(kT) + \int_{kT}^{(k+1)T} e^{A((k+1)T - \tau)} B u(\tau) \, d\tau
\end{equation*}
Dengan substitusi variabel $\lambda = (k+1)T - \tau$, diperoleh relasi matriks diskrit:
\begin{equation*}
    G = e^{AT}, \quad H = \left( \int_{0}^{T} e^{A\lambda} \, d\lambda \right) B
\end{equation*}
Jika matriks $A$ nonsingular (memiliki invers), rumus untuk $H$ dapat disederhanakan menjadi:
\begin{equation*}
    H = A^{-1}(e^{AT} - I) B
\end{equation*}

\subsection{Konversi Ruang Keadaan Diskrit ke Fungsi Transfer}
Dengan menerapkan Transformasi Z pada persamaan keadaan diskrit dengan $x(0) = 0$:
\begin{align*}
    z X(z) &= G X(z) + H U(z) \implies (zI - G) X(z) = H U(z) \\
    X(z) &= (zI - G)^{-1} H U(z)
\end{align*}
Substitusi ke persamaan keluaran $Y(z) = C X(z) + D U(z)$:

\kotakdefinisi{Fungsi Pulsa Transfer dari Ruang Keadaan Diskrit:}{
\begin{equation*}
    G(z) = \frac{Y(z)}{U(z)} = C (zI - G)^{-1} H + D = \frac{C \operatorname{adj}(zI - G) H}{\det(zI - G)} + D
\end{equation*}
}
Persamaan karakteristik sistem diskrit adalah $\det(zI - G) = 0$. Nilai-nilai eigen dari matriks $G$ ($\lambda_i$) merupakan kutub-kutub (\textit{poles}) sistem diskrit.

\vspace{6pt}

\section{Solusi Ruang Keadaan Diskrit \& Matriks Transisi Keadaan}

\subsection{Solusi Domain Waktu Rekursif}
Persamaan keadaan $x(k+1) = G x(k) + H u(k)$ dengan syarat awal $x(0)$ dapat diselesaikan secara rekursif:
\begin{align*}
    x(1) &= G x(0) + H u(0) \\
    x(2) &= G x(1) + H u(1) = G^2 x(0) + G H u(0) + H u(1) \\
    &\;\;\vdots \\
    x(k) &= G^k x(0) + \sum_{j=0}^{k-1} G^{k-1-j} H u(j)
\end{align*}
Dan persamaan keluaran:
\begin{equation*}
    y(k) = C G^k x(0) + C \sum_{j=0}^{k-1} G^{k-1-j} H u(j) + D u(k)
\end{equation*}

\subsection{Matriks Transisi Keadaan Diskrit \texorpdfstring{$\Psi(k) = G^k$}{Psi(k) = G^k}}
\kotakdefinisi{Perhitungan Matriks Transisi Keadaan Diskrit:}{
\begin{equation*}
    \Psi(k) = G^k = \mathcal{Z}^{-1}\left\{ (zI - G)^{-1} z \right\}
\end{equation*}
}

\textbf{Sifat-sifat Matriks Transisi Diskrit $G^k$}:
\begin{itemize}
    \item $G^0 = I$
    \item $G^{k_1 + k_2} = G^{k_1} G^{k_2} = G^{k_2} G^{k_1}$
    \item $(G^k)^{-1} = G^{-k}$ (jika $G$ nonsingular)
    \item $\Psi(k+1) = G \Psi(k) = \Psi(k) G$
\end{itemize}

\vspace{6pt}

\section{Kestabilan Sistem Waktu Diskrit}

\subsection{Pemetaan Bidang \texorpdfstring{$s$}{s} ke Bidang \texorpdfstring{$z$}{z}}
Hubungan $z = e^{sT}$ dengan $s = \sigma + j\omega$ menghasilkan:
\begin{equation*}
    z = e^{(\sigma + j\omega)T} = e^{\sigma T} e^{j\omega T} \implies |z| = e^{\sigma T}, \quad \angle z = \omega T
\end{equation*}

\begin{center}
\small
\begin{tabular}{|p{4.5cm}|p{4.5cm}|p{5.5cm}|}
\hline
\textbf{Daerah di Bidang $s$} & \textbf{Daerah di Bidang $z$} & \textbf{Karakteristik Kestabilan} \\ \hline
Setengah Bidang Kiri ($\sigma < 0$) & Bagian dalam lingkaran satuan ($|z| < 1$) & \textbf{Stabil Asimtotik} (Amplitudo meluruh) \\ \hline
Sumbu Imajiner ($\sigma = 0$) & Pada lingkaran satuan ($|z| = 1$) & \textbf{Stabil Marjinal} (Osilasi kontinu) \\ \hline
Setengah Bidang Kanan ($\sigma > 0$) & Bagian luar lingkaran satuan ($|z| > 1$) & \textbf{Tak Stabil} (Amplitudo divergen) \\ \hline
\end{tabular}
\end{center}

\kotakdefinisi{Kriteria Kestabilan Sistem Waktu Diskrit:}{
Sistem LTI diskrit bersifat \textbf{stabil asimtotik} jika dan hanya jika seluruh kutub dari fungsi transfer $G(z)$ (atau seluruh nilai eigen dari matriks sistem $G$) terletak \textbf{di dalam lingkaran satuan} pada bidang kompleks $z$:
\begin{equation*}
    |p_i| < 1 \quad \text{atau} \quad |\lambda_i| < 1, \quad \forall i = 1, 2, \dots, n
\end{equation*}
}

\subsection{Uji Kestabilan Jury (\textit{Jury Stability Test})}
Untuk polinomial karakteristik sistem orde-$n$:
\begin{equation*}
    P(z) = a_0 z^n + a_1 z^{n-1} + \dots + a_{n-1} z + a_n = 0 \quad (a_0 > 0)
\end{equation*}
Syarat perlu kestabilan:
\begin{enumerate}
    \item $P(1) > 0$
    \item $(-1)^n P(-1) > 0$
    \item $|a_n| < a_0$
\end{enumerate}
Untuk orde-2 ($P(z) = a_0 z^2 + a_1 z + a_2 = 0, a_0 > 0$), syarat perlu dan cukup adalah:
\begin{enumerate}
    \item $|a_2| < a_0$
    \item $P(1) = a_0 + a_1 + a_2 > 0$
    \item $P(-1) = a_0 - a_1 + a_2 > 0$
\end{enumerate}

\vspace{6pt}

\section{Contoh Soal dan Pembahasan Komprehensif}

\noindent\textbf{Contoh 1: Invers Transformasi Z}\\
Tentukan invers Transformasi Z dari fungsi:
\begin{equation*}
    X(z) = \frac{z(2z + 1)}{(z - 1)(z - 0.5)}
\end{equation*}

\textit{Penyelesaian:}

Bagi $X(z)$ dengan $z$ terlebih dahulu:
\begin{equation*}
    \frac{X(z)}{z} = \frac{2z + 1}{(z - 1)(z - 0.5)} = \frac{A}{z - 1} + \frac{B}{z - 0.5}
\end{equation*}
Menghitung koefisien $A$ dan $B$:
\begin{align*}
    A &= \lim_{z \to 1} (z - 1) \frac{X(z)}{z} = \frac{2(1) + 1}{1 - 0.5} = \frac{3}{0.5} = 6 \\
    B &= \lim_{z \to 0.5} (z - 0.5) \frac{X(z)}{z} = \frac{2(0.5) + 1}{0.5 - 1} = \frac{2}{-0.5} = -4
\end{align*}
Sehingga diperoleh:
\begin{equation*}
    X(z) = 6 \frac{z}{z - 1} - 4 \frac{z}{z - 0.5}
\end{equation*}
Dengan mengambil invers Transformasi Z dari masing-masing suku:
\begin{equation*}
    x(k) = 6(1)^k - 4(0.5)^k = 6 - 4(0.5)^k, \quad k = 0, 1, 2, \dots
\end{equation*}

\vspace{10pt}

\noindent\textbf{Contoh 2: Diskritisasi Plant Kontinu dengan ZOH}\\
Diberikan plant kontinu $G_p(s) = \dfrac{2}{s + 2}$ dengan periode pencuplikan $T = 0.1$ detik. Tentukan fungsi transfer diskrit $G(z)$.

\textit{Penyelesaian:}

Gunakan rumus diskritisasi dengan ZOH:
\begin{equation*}
    G(z) = (1 - z^{-1}) \mathcal{Z}\left\{ \frac{G_p(s)}{s} \right\} = \frac{z - 1}{z} \mathcal{Z}\left\{ \frac{2}{s(s+2)} \right\}
\end{equation*}
Ekspansi pecahan parsial pada domain $s$:
\begin{equation*}
    \frac{2}{s(s+2)} = \frac{1}{s} - \frac{1}{s+2}
\end{equation*}
Invers Laplace:
\begin{equation*}
    \mathcal{L}^{-1}\left\{ \frac{1}{s} - \frac{1}{s+2} \right\} = 1(t) - e^{-2t}
\end{equation*}
Lakukan pencuplikan pada $t = kT$:
\begin{equation*}
    x(kT) = 1 - e^{-2kT} = 1 - (e^{-2T})^k
\end{equation*}
Transformasi Z dari sinyal tercuplik:
\begin{equation*}
    \mathcal{Z}\left\{ 1 - (e^{-2T})^k \right\} = \frac{z}{z - 1} - \frac{z}{z - e^{-2T}} = \frac{z(1 - e^{-2T})}{(z - 1)(z - e^{-2T})}
\end{equation*}
Maka fungsi transfer diskrit $G(z)$:
\begin{equation*}
    G(z) = \frac{z - 1}{z} \cdot \frac{z(1 - e^{-2T})}{(z - 1)(z - e^{-2T})} = \frac{1 - e^{-2T}}{z - e^{-2T}}
\end{equation*}
Substitusikan $T = 0.1 \implies e^{-2(0.1)} = e^{-0.2} \approx 0.8187$:
\begin{equation*}
    G(z) = \frac{1 - 0.8187}{z - 0.8187} = \frac{0.1813}{z - 0.8187}
\end{equation*}

\vspace{10pt}

\noindent\textbf{Contoh 3: Ruang Keadaan Diskrit \& Matriks Transisi}\\
Diberikan sistem diskrit:
\begin{equation*}
    x(k+1) = \begin{bmatrix} 0 & 1 \\ -0.16 & 1 \end{bmatrix} x(k) + \begin{bmatrix} 0 \\ 1 \end{bmatrix} u(k), \quad y(k) = \begin{bmatrix} 1 & 0 \end{bmatrix} x(k)
\end{equation*}
Tentukan matriks transisi keadaan $G^k$ dan fungsi pulsa transfer $G(z)$.

\textit{Penyelesaian:}

Bentuk matriks $(zI - G)$:
\begin{equation*}
    zI - G = \begin{bmatrix} z & -1 \\ 0.16 & z - 1 \end{bmatrix}
\end{equation*}
Determinan:
\begin{equation*}
    \det(zI - G) = z(z - 1) - (-1)(0.16) = z^2 - z + 0.16 = (z - 0.8)(z - 0.2)
\end{equation*}
Invers $(zI - G)^{-1}$:
\begin{equation*}
    (zI - G)^{-1} = \frac{1}{(z - 0.8)(z - 0.2)} \begin{bmatrix} z - 1 & 1 \\ -0.16 & z \end{bmatrix}
\end{equation*}

Menghitung fungsi transfer $G(z) = C (zI - G)^{-1} H$:
\begin{align*}
    G(z) &= \begin{bmatrix} 1 & 0 \end{bmatrix} \left( \frac{1}{(z - 0.8)(z - 0.2)} \begin{bmatrix} z - 1 & 1 \\ -0.16 & z \end{bmatrix} \right) \begin{bmatrix} 0 \\ 1 \end{bmatrix} \\
    &= \frac{1}{(z - 0.8)(z - 0.2)} \begin{bmatrix} 1 & 0 \end{bmatrix} \begin{bmatrix} 1 \\ z \end{bmatrix} = \frac{1}{z^2 - z + 0.16}
\end{align*}

Menghitung matriks transisi keadaan $G^k = \mathcal{Z}^{-1}\{(zI - G)^{-1} z\}$:
\begin{equation*}
    (zI - G)^{-1} z = \frac{z}{(z - 0.8)(z - 0.2)} \begin{bmatrix} z - 1 & 1 \\ -0.16 & z \end{bmatrix}
\end{equation*}
Dekomposisi pecahan parsial tiap elemen:
\begin{align*}
    \frac{(z - 1)z}{(z - 0.8)(z - 0.2)} &= -\frac{1}{3}\frac{z}{z - 0.8} - \frac{4}{3}\frac{z}{z - 0.2} \implies -\frac{1}{3}(0.8)^k - \frac{4}{3}(0.2)^k \\
    \frac{z}{(z - 0.8)(z - 0.2)} &= \frac{5}{3}\frac{z}{z - 0.8} - \frac{5}{3}\frac{z}{z - 0.2} \implies \frac{5}{3}(0.8)^k - \frac{5}{3}(0.2)^k \\
    \frac{-0.16 z}{(z - 0.8)(z - 0.2)} &= -\frac{4}{15}\frac{z}{z - 0.8} + \frac{4}{15}\frac{z}{z - 0.2} \implies -\frac{4}{15}(0.8)^k + \frac{4}{15}(0.2)^k \\
    \frac{z^2}{(z - 0.8)(z - 0.2)} &= \frac{4}{3}\frac{z}{z - 0.8} - \frac{1}{3}\frac{z}{z - 0.2} \implies \frac{4}{3}(0.8)^k - \frac{1}{3}(0.2)^k
\end{align*}
Sehingga matriks transisi keadaan diperoleh:
\begin{equation*}
    G^k = \begin{bmatrix}
        -\frac{1}{3}(0.8)^k - \frac{4}{3}(0.2)^k & \frac{5}{3}(0.8)^k - \frac{5}{3}(0.2)^k \\[6pt]
        -\frac{4}{15}(0.8)^k + \frac{4}{15}(0.2)^k & \frac{4}{3}(0.8)^k - \frac{1}{3}(0.2)^k
    \end{bmatrix}, \quad k \ge 0
\end{equation*}

\end{document}
